% Image credits for the photographs used in this volume (Book 1).
% Machine-readable license records live in images/book1/CREDITS.md; keep
% the two files in sync when adding or removing a photograph.
\chapter*{Image Credits}

The drawings in this book are original. The photographs are used
under their respective free licenses, with thanks to their authors:

\begin{itemize}
  \item \emph{John Dalton} (atoms and molecules chapter): engraving by
        Charles Turner after Joseph Allen, 1834, Library of Congress via
        Wikimedia Commons, public domain.
  \item \emph{Dalton's table of the elements, 1808} (atoms and molecules
        chapter): plate from John Dalton's \emph{A New System of Chemical
        Philosophy}, scan by the Internet Archive via Wikimedia Commons,
        public domain.
  \item \emph{Bauxite} (raw materials chapter): photograph by Werner
        Schellmann via Wikimedia Commons, CC BY-SA 2.5.
  \item \emph{Granite} (pure substances chapter): photograph by Eurico
        Zimbres via Wikimedia Commons, CC BY-SA 2.0 br.
  \item \emph{Anhydrous copper sulfate} (identifying substances chapter):
        photograph by W.~Oelen via Wikimedia Commons, CC BY-SA 3.0.
  \item \emph{Copper sulfate crystals} (identifying substances chapter):
        photograph by Stephanb via Wikimedia Commons, CC BY-SA 3.0.
  \item \emph{Antoine-Laurent Lavoisier and Marie-Anne Paulze Lavoisier}
        (balanced equations chapter): painting by Jacques-Louis David,
        1788, Metropolitan Museum of Art via Wikimedia Commons, public
        domain.
  \item \emph{Bakelite telephone} (plastics chapter): photograph by Holger
        Ellgaard via Wikimedia Commons, CC BY-SA 3.0.
  \item \emph{Ernest Rutherford} (inside the atom chapter): photograph,
        George Grantham Bain Collection, Library of Congress, via
        Wikimedia Commons, public domain.
  \item \emph{Halite} (ions chapter): photograph by Robert M.~Lavinsky
        via Wikimedia Commons, CC BY-SA 3.0.
  \item \emph{The Statue of Liberty} (metals and corrosion chapter):
        photograph by Elcobbola via Wikimedia Commons, public domain.
  \item \emph{Dmitri Mendeleev} and \emph{his table of 1869} (periodic
        table chapter): via Wikimedia Commons, public domain.
  \item \emph{The Crab Nebula} (elements chapter): NASA, ESA, J.~Hester
        and A.~Loll, public domain.
  \item \emph{Sodium} (electron shells chapter): photograph by Dnn87 via
        Wikimedia Commons, CC BY-SA 3.0.
  \item \emph{Chlorine in an ampoule} (electron shells chapter):
        photograph by W.~Oelen via Wikimedia Commons, CC BY-SA 3.0.
  \item \emph{Amedeo Avogadro} (mole chapter): lithograph after a
        drawing by C.~Sentier, 1856, Edgar Fahs Smith collection, via
        Wikimedia Commons, public domain.
  \item \emph{Snow crystals} (polarity chapter): photographs by Wilson
        Bentley, 1902, via Wikimedia Commons, public domain.
  \item \emph{Louis Pasteur} (stereochemistry chapter): photograph by
        Paul Nadar, via Wikimedia Commons, public domain.
  \item \emph{Ammonium tartrate crystals} (stereochemistry chapter):
        photograph by Marie-Lan Taÿ Pamart via Wikimedia Commons,
        CC BY 4.0.
  \item \emph{Svante Arrhenius} (acids and bases chapter): photogravure,
        1909, via Wikimedia Commons, public domain.
  \item \emph{A voltaic pile} (cells chapter): photograph by GuidoB via
        Wikimedia Commons, CC BY-SA 3.0.
  \item \emph{Wallace Carothers in his laboratory} (polymers chapter):
        photographer unknown, via Wikimedia Commons, public domain.
\end{itemize}

Full source links and license URLs for every photograph are listed
in \texttt{images/book1/CREDITS.md} in the book's repository.
The hazard pictograms are those of the Globally Harmonized System of
Classification and Labelling of Chemicals, as drawn for the
\texttt{ghsystem} package of \TeX\ Live.

\bigskip
Alongside the photographs and the original drawings, some figures are
AI-generated illustrations, produced with OpenAI image generation from
prompts written by the book's editors, and reviewed for chemical
accuracy before inclusion. The full prompt for every generated image is
recorded in \texttt{images/book1/ai/PROMPTS.md} in the repository.
\clearpage
